\section*{Capítulo \ref{ch:b3:inorganic-materials} --- Materiais inorgânicos e nanomateriais}

\begin{solution}{exo:b3:inorganic-materials:1}
$x^2 \propto t$: uma espessura três vezes maior leva nove vezes mais tempo, \qty{90}{h}.
\end{solution}

\begin{solution}{exo:b3:inorganic-materials:2}
Hidrólise: $\ce{Si(OC2H5)4 + 4H2O -> Si(OH)4 + 4C2H5OH}$.\\
Condensação: $\ce{2Si(OH)4 -> Si2O(OH)6 + H2O}$.\\
Global: $\ce{Si(OC2H5)4 + 2H2O -> SiO2 + 4C2H5OH}$.
\end{solution}

\begin{solution}{exo:b3:inorganic-materials:3}
Formadores: $\ce{SiO2}$, $\ce{B2O3}$. Modificadores: $\ce{Na2O}$, $\ce{CaO}$. Cada íon óxido do $\ce{Na2O}$ rompe uma ponte Si--O--Si
em duas extremidades Si--O$^-$, cada uma acompanhada de um $\ce{Na+}$: a rede é cortada, o fundido
escoa mais facilmente e funde a temperatura mais baixa.
\end{solution}

\begin{solution}{exo:b3:inorganic-materials:4}
Doze pentágonos (como em todo \omterm{def:b3:inorganic-materials:carbon}{fulereno}). $V = 70$, $E = 3V/2 = 105$ ligações, $F = 2 - V + E = 37$
faces, logo $37 - 12 = 25$ hexágonos.
\end{solution}

\begin{solution}{exo:b3:inorganic-materials:5}
$\sqrt2\,(r_B + r_O) = 1.4142 \times \qty{200.5}{pm} = \qty{283.5}{pm}$. $\ce{SrTiO3}$: $284/283.5 = 1.00$, a perovskita cúbica ideal.
$\ce{BaTiO3}$: $301/283.5 = 1.06$, titânio pequeno demais para seu octaedro, fora do centro: \omterm{def:b3:inorganic-materials:ferroelectric}{ferroelétrico}.
$\ce{CaTiO3}$: $274/283.5 = 0.97$, cálcio pequeno demais para sua cavidade: octaedros inclinados, simetria
mais baixa.
\end{solution}

\begin{solution}{exo:b3:inorganic-materials:6}
Si/Al $= 106/86 = 1.23$. $M = 86 \times 23.0 + 86 \times 27.0 + 106 \times 28.1 + 384 \times 16.0 = \qty{13422.6}{g/mol}$; 43 $\ce{Ca^2+}$ por
fórmula: $43/\qty{13422.6}{g/mol} = \qty{3.20}{mmol/g}$.
\end{solution}

\begin{solution}{exo:b3:inorganic-materials:7}
Distância entre vizinhos mais próximos $d = a/\sqrt2 = \qty{288.4}{pm}$. Um aglomerado de $n$ camadas mede $(2n + 1)d$ de diâmetro:
$2n + 1 \approx 5/0.2884 = 17.3$, logo $n = 8$ (\qty{4.9}{nm}). $N(8) = 2057$ átomos, $10 \times 64 + 2 = 642$ na superfície:
31\,\%.
\end{solution}

\begin{solution}{exo:b3:inorganic-materials:8}
$h^2/(8\mu R^2)$ com $\mu = \qty{9.109e-32}{kg}$: \qty{0.94}{eV} a \qty{2.0}{nm} e \qty{0.42}{eV} a \qty{3.0}{nm}. Emissão em
$1.74 + 0.94 = \qty{2.68}{eV}$, \qty{463}{nm} (azul), e $1.74 + 0.42 = \qty{2.16}{eV}$, \qty{574}{nm} (amarelo): o ponto menor
emite a luz mais azul.
\end{solution}

\begin{solution}{exo:b3:inorganic-materials:9}
$(10, 10)$: poltrona, $n - m = 0$, metálico. $(10, 0)$: zigue-zague, 10 não é múltiplo de 3, \omterm{def:b3:solid-state:classes}{semicondutor}.
$(9, 0)$: zigue-zague, metálico. $(12, 8)$: nenhum dos dois (quiral), $n - m = 4$, \omterm{def:b3:solid-state:classes}{semicondutor}.
\end{solution}

\begin{solution}{exo:b3:inorganic-materials:10}
$M = 4 \times 65.4 + 13 \times 16.0 + 24 \times 12.0 + 12 \times 1.0 = \qty{769.6}{g/mol}$; $a^3 = (\qty{25.82e-8}{cm})^3 = \qty{1.721e-20}{cm^3}$;
$\rho = 8 \times 769.6/(\num{6.022e23} \times \num{1.721e-20}) = \qty{0.594}{g.cm^{-3}}$. Um campo mede \qty{7140}{m^2}: um grama corresponde a
$3000/7140 = 0.42$ campo, e uma colher de chá de alguns gramas, a mais que um campo inteiro.
\end{solution}

\begin{solution}{exo:b3:inorganic-materials:11}
$\dfrac{10n^2 + 2}{(10n^3 + 15n^2 + 11n + 3)/3} = \dfrac{30n^2 + 6}{10n^3 + 15n^2 + 11n + 3} \to \dfrac{30n^2}{10n^3} = \dfrac3n$. Para $n = 8$: exata
$642/2057 = 0.31$, contra $3/8 = 0.375$; a aproximação superestima para os aglomerados
pequenos, cujas camadas internas não são desprezíveis.
\end{solution}

\begin{solution}{exo:b3:inorganic-materials:12}
$2E = 3V$, $2E = 5p + 6h + 7s$, $F = p + h + s$. Euler: $V - E + F = 2$, com $V = 2E/3$: $F - E/3 = 2$, logo
$p + h + s - (5p + 6h + 7s)/6 = 2$, isto é, $(p - s)/6 = 2$: $p - s = 12$. Cada heptágono exige um pentágono
a mais.
\end{solution}

\begin{solution}{pb:b3:inorganic-materials:1}
\textbf{1.} $\mathrm{Na_{12}Al_{12}Si_{12}O_{48}}$: $12 \times 23.0 + 12 \times 27.0 + 12 \times 28.1 + 48 \times 16.0 = \qty{1705.2}{g/mol}$.

\textbf{2.} $1705.2 + 27 \times 18.0 = \qty{2191.2}{g/mol}$.

\textbf{3.} Si/Al $= 12/12 = 1$.

\textbf{4.} Nenhuma ligação Al--O--Al; com Si/Al $= 1$, todo Si tem quatro vizinhos Al e todo
Al, quatro Si: alternância estrita.

\textbf{5.} $-12$, uma por alumínio.

\textbf{6.} $486.0/2191.2 = 22.2$\,\%.

\textbf{7.} $\ce{2Na+}(\text{zeólita}) + \ce{Ca^2+}(\text{aq}) \rightleftharpoons \ce{Ca^2+}(\text{zeólita}) + \ce{2Na+}(\text{aq})$.

\textbf{8.} Seis (doze cargas).

\textbf{9.} $6/\qty{1705.2}{g/mol} = \qty{3.52}{mmol/g}$.

\textbf{10.} $6/\qty{2191.2}{g/mol} = \qty{2.74}{mmol/g}$.

\textbf{11.} $\qty{3.52}{mmol/g} \times \qty{100.1}{g/mol} = \qty{352}{mg}$ de $\ce{CaCO3}$ por grama.

\textbf{12.} $\qty{2.0}{mmol/L} \times \qty{15}{L} = \qty{30}{mmol}$.

\textbf{13.} $\qty{30}{mmol}/\qty{2.74}{mmol/g} = \qty{11}{g}$.

\textbf{14.} $\qty{10.9}{g}/0.60 = \qty{18}{g}$.

\textbf{15.} $\qty{60}{mmol}$ de $\ce{Na+}$, $\qty{0.060}{mol} \times \qty{23.0}{g/mol} = \qty{1.4}{g}$.

\textbf{16.} Os fosfatos das águas residuais alimentam as algas nos rios e lagos; as florações
e sua decomposição consomem o oxigênio dissolvido (eutrofização).

\textbf{17.} O $\ce{Mg^2+}$, menor, retém suas moléculas de água mais fortemente; ele precisa perdê-las
para passar pela janela e chegar aos sítios, o que é lento nas temperaturas de lavagem.

\textbf{18.} O íon nu mede cerca de \qty{0.20}{nm}, metade da janela.

\textbf{19.} O íon hidratado é maior que a janela: ele precisa perder parte de sua
camada de hidratação para passar, e os oxigênios da estrutura tomam o lugar da
água.

\textbf{20.} A estrutura tem carga negativa e só troca cátions; um
ânion grande de cadeia longa é repelido e não pode passar pela janela, de modo que fica
na água, onde faz seu trabalho.

\textbf{21.} Desidratada, a \omterm{def:b3:inorganic-materials:zeolite}{zeólita} capta fortemente a água em suas cavidades; só
as moléculas menores que a janela entram, de modo que ela separa as moléculas pelo tamanho:
uma peneira em escala molecular.

\textbf{22.} Os íons $\ce{K+}$, maiores, ficam nas janelas e as estreitam: só moléculas
menores, como a água, são admitidas.

\textbf{23.} O isômero \textit{para}, esguio e linear, cabe nos canais e se difunde depressa; os
isômeros \textit{orto} e \textit{meta}, mais largos, se difundem lentamente, ficam e se isomerizam.

\textbf{24.} A \omterm{def:b3:inorganic-materials:zeolite}{zeólita} A anidra pode captar, em teoria, \qty{3.52}{mmol} de $\ce{Ca^2+}$ por grama.
\end{solution}
